\section{The GIJS Distinguisher In the Our Context}\label{sec:gijs}
In this section we analyze the performance of the GIJS
distinguisher~\cite{GIJS26} in our attack algorithms.

We now show that under plausible and empirically validated heuristics, the GIJS
distinguisher algorithms succeeds in distinguishing
clean columns (shortened from a pure GRS column) from dirty columns (one of two
elements of a twin pair, both of which appear in the shortened code) under
certain
conditions even with noisy RLCE columns.

The shortened clean columns lie on a projected
rational curve
\[
    \varphi : \mathbf{P}^1 \longrightarrow \mathbf{P}^{k'-1},
\]
of degree $D=k'+w-1$, given by homogeneous polynomials of degree $D$ in the
decoding
parameter $t$. Below, we work over the algebraic closure
$\overline{\mathbf{F}}_q$ when counting roots.

Each original column has a distinct evaluation parameter
$\alpha \in \mathbf{F}_q$, and this distinct parameter is unchanged by
shortening.

Fix a holdout column and let $I$ be the set of interpolation
columns, $|I| = n'-1$. Partition $I$ as follows:
\begin{itemize}
    \item \emph{Clean points:} kept pure columns and singleton-hit
          survivors. Each lies on the curve $\varphi(\mathbf{P}^1)$,
          at the parameter of its original column. Let $\mathcal{C}$
          denote this set, and $C = |\mathcal{C}|$.
    \item \emph{Dirty points:} kept columns belonging to a pair both
          of whose members are kept. Such a pair is called
          \emph{both-kept}. Let $\mathcal{U}$ denote the set of
          non-degenerate both-kept pairs, and $U = |\mathcal{U}|$.
          Non-degenerate means that neither mixing coefficient is zero,
          so neither member of the pair is itself a curve point; in the
          RLCE specification the mixing coefficients are chosen to always
          be non-degenerate.
\end{itemize}
Thus, counting kept columns, $n' = C + 2U$ plus any degenerate
both-kept pairs, which we exclude from $\mathcal{U}$ and treat
separately; in the ID-0 run below there are none.

Let $H$ be a homogeneous polynomial of degree $d$ in $k'$
variables. We say $H$ \emph{vanishes to order $s$} at a point
$x$ if all partial derivatives of $H$ of order $< s$ vanish at $x$,
i.e.\ the jet of $H$ at $x$ of order $s-1$ is zero. The GIJS
kernel, for a fixed interpolation set, is the space of such $H$
vanishing to order $s$ at every point of $I$. A clean holdout
$\tau$ is accepted, i.e.\ not flagged, if every $H$ in the kernel
vanishes at $\tau$; it is flagged if some $H$ in the kernel does
not vanish at $\tau$.

\subsection{PIR persistence: rigorous acceptance with
    noise}\label{sec:pir-noise}

First, we show that the pairs of dirty columns (which are the only type of
non-clean columns present in the shortened code)
are actually useful to the GIJS attack. As we will show, the reason is
precisely due to the linear relation existing between pairs that results in
a pure GRS column when inverting the $A$ map from public key generation.

\begin{theorem}[PIR persistence]
    Let $H$ be a homogeneous polynomial of degree $d$ vanishing to
    order $s$ at every interpolation point, and assume
    \begin{equation}
        2s > d.
        \label{eq:line}
    \end{equation}
    Then $H$ vanishes identically on the line through the two dirty
    columns of every non-degenerate both-kept pair, and in particular
    $H$ vanishes at the underlying GRS point of that pair.
    Consequently, if
    \begin{equation}
        sC + U > dD,
        \label{eq:condition}
    \end{equation}
    then $H$ vanishes identically on the curve $\varphi(\mathbf{P}^1)$.

    As a result, it vanishes at every clean holdout column, meaning that
    running
    the
    GIJS distinguisher with a clean holdout column outputs "GRS"
    deterministically
    if the above conditions are fulfilled.
\end{theorem}

\begin{proof}
    Fix a non-degenerate both-kept pair with underlying GRS point
    $P$ and random point $R$. After shortening, its two columns are
    \[
        Q_1 = aP + cR, \qquad Q_2 = bP + dR,
    \]
    in projective coordinates, with $ad - bc \ne 0$ and $c, d \ne 0$
    by non-degeneracy. Hence $Q_1 \ne Q_2$ in $\mathbf{P}^{k'-1}$,
    and $Q_1, Q_2, P$ are collinear on the line $L$ spanned by $P$
    and $R$; $L$ is not contained in the curve, since $R$ is not a
    curve point.

    Restrict $H$ to $L$. The restriction is a binary form of degree
    $d$ on $L$, vanishing to order $s$ at the two distinct points
    $Q_1$ and $Q_2$. A binary form of degree $d$ with two distinct
    roots of multiplicity $s$ each is identically zero as soon as
    $2s > d$, by degree counting. This is~\eqref{eq:line}. Hence
    $H$ vanishes identically on $L$, and in particular $H(P) = 0$.

    Now consider the pullback $h(t) = H(\varphi(t))$. This is a binary
    form of degree $dD$, unless it is identically zero, which is the
    conclusion we want. We count its roots over
    $\overline{\mathbf{F}}_q$, with multiplicity.

    First, let $t_0$ be the parameter of a clean interpolation point
    $X = \varphi(t_0)$. Since $H$ vanishes to order $s$ at $X$ in
    ambient space, its Taylor expansion at $X$ has no terms of degree
    $< s$. Composing with $\varphi(t) - X$, which has zero constant
    term at $t_0$, gives a power series in $t - t_0$ with no terms of
    degree $< s$; if $\varphi$ is ramified at $t_0$, the order is only
    larger. Hence $h$ vanishes to order at least $s$ at $t_0$.

    Second, let $t_i$ be the parameter of the underlying GRS point
    $P_i$ of a non-degenerate both-kept pair. By the line argument
    above, $H(P_i) = 0$, so $h(t_i) = 0$, i.e.\ $h$ has a root at
    $t_i$ of order at least $1$. The parameters $t_i$ are distinct
    across pairs, since the original evaluation parameters are
    distinct, and they are distinct from the clean parameters, since
    a both-kept pair's parameter is not the parameter of any clean
    interpolation column. Hence the roots counted in the first and
    second steps are distinct, except that a clean root already
    carries multiplicity $s$.

    Therefore $h$ has total root multiplicity at least $sC + U$.
    If \eqref{eq:condition} holds, this exceeds the degree $dD$ of
    $h$, so $h$ cannot be a nonzero binary form. Hence $h \equiv 0$,
    i.e.\ $H$ vanishes at every point of the curve
    $\varphi(\mathbf{P}^1)$. A clean holdout column is a curve
    point, so every kernel polynomial $H$ vanishes at it, and the
    GIJS test accepts it.
\end{proof}

\subsection{Random Holdout Columns}\label{sec:rand-holdout}
We now show that with high probability (and heuristically with overwhelming
probability) that a random and independent column
will be flagged by the GIJS distinguisher, meaning that at least one kernel
polynomial $H$ will not vanish at it.

\begin{theorem}[Independent uniform random holdout]
    Fix the interpolation set $I$, and let $K$ be its GIJS kernel
    as above. Assume $K \ne \{0\}$. Let $\tau$ be uniform in
    $\mathbf{F}_q^{k'}$, independent of $I$ and hence independent
    of $K$. Then
    \begin{equation}
        \Pr[\tau \text{ is flagged}] \ge 1 - \frac{d}{q}.
        \label{eq:bound}
    \end{equation}
\end{theorem}

\begin{proof}
    Since $K \ne \{0\}$, pick any nonzero $H_0 \in K$. The flagging
    event contains the event $\{H_0(\tau) \ne 0\}$, because if
    $H_0$ does not vanish at $\tau$, the GIJS test has a witness
    and flags $\tau$. Hence
    \[
        \Pr[\tau \text{ is flagged}]
        \ge \Pr[H_0(\tau) \ne 0]
        = 1 - \Pr[H_0(\tau) = 0].
    \]
    It remains to bound $\Pr[H_0(\tau) = 0]$. This is the
    Schwartz--Zippel bound: a nonzero polynomial of total degree
    $d$ in $k'$ variables over $\mathbf{F}_q$ vanishes at at most
    a $d/q$ fraction of points in $\mathbf{F}_q^{k'}$. Indeed, the
    bound follows by induction on $k'$, writing $H_0$ as a
    polynomial in one variable with coefficients in the remaining
    variables; a univariate polynomial of degree at most $d$ has at
    most $d$ roots unless it is identically zero. Applying it to
    $H_0$ gives $\Pr[H_0(\tau) = 0] \le d/q$, and
    \eqref{eq:bound} follows.
\end{proof}

\subsubsection{Non-zero Kernel Conditions}\label{sec:nonzero-kernel}
It remains to discuss conditions on which the kernel is non-zero. 

Let $J=\binom{k'+s-1}{s-1}$ be the number of jet conditions, and 
$N_d=\binom{k'+d-1}{d}$ the number of monomials of degree $d$ in $k'$ variables.

\begin{corollary}[Sufficient condition for a nonzero kernel]
    If
    \begin{equation}
        (n'-1)J < N_d,
        \label{eq:kernel}
    \end{equation}
    then $K \ne \{0\}$ for every interpolation set, and hence
    \eqref{eq:bound} holds for every trial.
\end{corollary}

\begin{proof}
    $K$ is the kernel of an $(n'-1)J \times N_d$ matrix. Its rank
    is at most the number of rows, $(n'-1)J$. By rank--nullity,
    \[
        \dim K \ge N_d - (n'-1)J > 0
    \]
    under \eqref{eq:kernel}, so $K \ne \{0\}$. Theorem~3 then
    applies.
\end{proof}


\subsubsection{Heuristic Bounds on Dirty Columns}\label{sec:heuristic}

As our experiments (see below) show, we have not had \emph{any} dirty columns 
fail to flag (i.e. the GIJS distinguisher outputs "GRS" when they are a holdout column). 

Under a reasonable heuristic and using the $N_d - (n'-1)J$ lower bound for the kernel dimension, it seems 
reasonable to assume that the probability of failure in practice is more on the order of $(d/q)^{\dim K}$, 
meaning the probability is essentially independent across kernel polynomials. We stress
that this is a pure heuristic, but that it has been experimentally shown to be reasonable on 
parameter sets up to the lowest set given in the RLCE specification. 
